\documentclass[../main.tex]{subfiles}

\begin{document}

% ===========================================================================
\subsection{Conclusions}
% ===========================================================================

This project set out to answer a deceptively simple question: can a medical database be simultaneously GDPR-compliant by design and performant enough for real-time clinical use? The work presented in Part~I provides a concrete, empirically validated answer: yes, provided the right architectural decisions are made at the schema level rather than delegated to application code or access-control layers. From there, the project grew to span three complementary paradigms ---relational, semi-structured and document-oriented---, which made it possible to verify first-hand that there is no single \enquote{right} technology, only the most suitable one for each kind of clinical data.

\paragraph{On privacy-by-design compliance.}
The SHA-256-based separation between \texttt{paciente\_anonimo} and \texttt{datos\_personales} achieves genuine structural pseudonymisation as defined in Recital~26 of the GDPR. The critical property is that the anonymisation is not a removable layer: the investigator user account is denied access to \texttt{datos\_personales} at the database engine level, not through application logic that could be bypassed. Linking a \texttt{codigo\_hash} back to an individual identity requires either the original NIF (which lives outside the database) or a brute-force preimage attack against SHA-256, currently considered computationally infeasible. The \texttt{ON DELETE RESTRICT} constraints further guarantee that no record can be deleted from the clinical model while its personal data counterpart still exists, preserving audit-trail integrity.

\paragraph{On multi-engine optimisation.}
The three storage engines chosen for ChronoHealthDB each delivered measurable benefits in their respective roles:

\begin{itemize}
    \item \textbf{InnoDB} provided the transactional guarantees, referential integrity and multi-version concurrency required for the clinical workflow tables. Query benchmarking confirmed that composite indexes (\texttt{idx\_paciente\_fecha}) reduced patient-history join times from 39.89~ms to 0.64~ms (a \textbf{62.3× improvement}), and temporal aggregation indexes (\texttt{idx\_fecha\_evento}) reduced monthly-count queries from 46.06~ms to 1.36~ms (\textbf{33.8×}).

    \item \textbf{ARCHIVE} demonstrated that storage optimisation for immutable historical records is not a marginal concern: the same dataset that occupies 2,560~KiB under InnoDB compresses to 533~KiB under ARCHIVE, a \textbf{79.2\% reduction} attributable to row-level zlib compression and the absence of MVCC overhead. At the scale of a regional health authority managing 200,000 medical leaves per year over a decade, this translates to a difference of over 640~GB of primary storage.

    \item \textbf{MEMORY} confirmed its suitability for real-time queue management: triage priority queries executed in 0.77~ms, \textbf{45.9\% faster} than the equivalent InnoDB table. More importantly, MEMORY guarantees constant latency independent of disk load, a property that InnoDB cannot provide when its buffer pool is under pressure.
\end{itemize}

\paragraph{On the data generation pipeline.}
The clinically coherent synthetic dataset of 40,000 patients (with scenario-driven text templates for clinical descriptions, vital signs correlated with pathology, and age-stratified risk factors) validated that the benchmarking results are not artefacts of trivially simple data. The dataset behaves like real-world data from the query-pattern perspective, lending credibility to the performance measurements.

\paragraph{On the semi-structured phase (Part~II).}
Moving the same data into XML and loading it into eXist-db showed that the hierarchical nature of a clinical record is better exploited when stored as a tree than when scattered across tables. The three XQuery queries worked as a genuine decision-support system: the triangulation crossed emergencies, history and sick-leave records in a single FLWOR expression; the chronic-patient detection was rewritten from a near-quadratic cost to a linear one by relying on the document structure itself; and the hypoxia monitor forced careful type handling (\texttt{castable as xs:double}), a far from trivial detail when working with data that is always text in XML. The takeaway is that the effort spent giving the documents a good shape pays off handsomely when querying them.

\paragraph{On the hybrid MongoDB architecture (Part~III).}
The NoSQL phase confirmed that MySQL and MongoDB can coexist by assigning each what it does best: the transactional layer to the relational engine and the flexible clinical knowledge (the guidelines) to the document store. Three points stood out. First, privacy-by-design is genuinely feasible: hashing the national ID with SHA-256 at insertion time meant the 50,000 records reached MongoDB already pseudonymised, with no exposure window. Second, \texttt{.explain()} analysis reproduces in a NoSQL engine the very same indexing lesson seen in MySQL: the \texttt{anio\_nacimiento\_1} index turned a \texttt{COLLSCAN} over 50,000 documents into an \texttt{IXSCAN} that examines only the 12,040 relevant ones. Third, delegating the heavy lifting to the engine (the aggregation pipelines) and keeping Python as a mere presentation layer is the architectural decision that separates a solution that scales from an anti-pattern that exhausts client memory.

\paragraph{Overall assessment.}
The project demonstrates that the tension between GDPR compliance and operational performance is a false dilemma. Compliance by design, through structural data separation, and performance optimisation, through engine selection and indexing, are not competing goals but complementary architectural decisions that reinforce each other. Beyond the relational layer, walking through the three paradigms made it clear that the choice of storage technology is not a matter of fashion but of matching the tool to the shape of the data: tables for the tabular, XML trees for the hierarchical and JSON documents for the flexible. The added complexity of handling three technologies is fully justified by what each one contributes where the others fall short.

% ===========================================================================
\subsection{Future lines of research}
% ===========================================================================

The relational phase presented here constitutes the first layer of a naturally extensible system. Several research directions emerge directly from its limitations:

\paragraph{Graph databases for patient relationships.}
Parts~II and~III already covered the semi-structured (XML) and document-oriented (MongoDB) paradigms, so the next natural leap is towards \textbf{graph databases} (Neo4j or Apache AGE on PostgreSQL). Clinical networks are inherently relational in the graph-theoretic sense: co-morbidity networks, drug-interaction graphs, care-pathway analyses and epidemic contact tracing are queries that express naturally as graph traversals but are expensive as multi-join SQL operations or even as document aggregations. A compelling extension would be to project the \texttt{evento\_clinico} table into a property graph and evaluate algorithms such as shortest path or community detection on the already pseudonymised data.

\paragraph{Free-text clinical content and semantic search.}
The current model structures events well, but free-text clinical notes, imaging reports and discharge summaries remain the major challenge. A natural direction would be to exploit MongoDB's full-text search capabilities ---or to add vector indexes and language models--- to perform semantic search over that unstructured content, always preserving the hash-based pseudonymisation that the rest of the system already guarantees.

\paragraph{Horizontal scalability and replication.}
The current system is a single-node MySQL instance. A production deployment at hospital scale would require a replication topology (InnoDB Cluster or MySQL Group Replication) with read replicas for the research workload and a primary node for clinical writes. The effect of replication lag on MEMORY tables —which are local to each node— is an open engineering question worth investigating.

\paragraph{Encrypted ARCHIVE tables.}
MySQL's InnoDB tablespace encryption (\texttt{ENCRYPTION='Y'}) is well documented, but ARCHIVE engine encryption is not natively supported in all configurations. Exploring transparent disk-level encryption for ARCHIVE tables, or migrating the historical layer to a TDE-capable engine (e.g.\ Percona's TokuDB with compression), would address the regulatory requirement to protect data at rest without sacrificing the compression benefits demonstrated in this work.

\paragraph{Machine learning on anonymised data.}
The 40,000-patient dataset already contains enough signal to train basic predictive models on pseudonymised data: readmission risk, length-of-stay estimation, or early warning scores based on triage vital signs. A privacy-preserving ML pipeline that operates exclusively on the \texttt{investigador}-accessible tables would demonstrate a concrete downstream application of the GDPR-compliant architecture.

\end{document}
